% Source: https://artofproblemsolving.com/community/c3h1592668_rationalizing_denominators
% Stable identity: aops-topic-1592668-q01
\textbf{(a)} Take out the square factors: $\sqrt{63x}=\sqrt{9\cdot7x}=3\sqrt{7x}$ and $\sqrt{8b^2}=\sqrt{4b^2\cdot2}=2b\sqrt2$ (using $b>0$). Then multiply the numerator and the denominator by $\sqrt2$:
\[
\dfrac{\sqrt{63x}}{\sqrt{8b^2}}=\dfrac{3\sqrt{7x}}{2b\sqrt2}=\dfrac{3\sqrt{7x}\cdot\sqrt2}{2b\cdot2}=\boxed{\dfrac{3\sqrt{14x}}{4b}}.
\]

\textbf{(b)} Multiplying by $\dfrac{\sqrt{8b^2}}{\sqrt{8b^2}}$ is a correct method, but the numerator becomes $\sqrt{63x}\cdot\sqrt{8b^2}=\sqrt{504xb^2}$, \emph{not} $504xb^2$: the student dropped the square root. Done correctly, $504=36\cdot14$, so $\sqrt{504xb^2}=6b\sqrt{14x}$ and
\[
\dfrac{\sqrt{63x}\cdot\sqrt{8b^2}}{8b^2}=\dfrac{6b\sqrt{14x}}{8b^2}=\dfrac{3\sqrt{14x}}{4b},
\]
the same answer. Multiplying by $\sqrt2$ instead is simply quicker, because $\sqrt2$ is the only irrational factor left in the simplified denominator.
